% Table: The role of degree choice (Chapter 4)
% NOTE: the source file referenced \ref{tab:earnings}, a label that does not
% exist in the paper; the linear IV estimate it points to is in tab:main.
\begin{table}[htbp]
\centering
\caption{The role of degree choice}
\label{tab:degree_channel}
\begin{tabular}{l ccc}
\toprule
                                      & Reduced      & IV           & IV             \\
                                      & form         & (linear)     & (saturated)    \\
                                      & (1)          & (2)          & (3)            \\
\midrule
\multicolumn{4}{c}{\textbf{Panel A: Effect on studying economics at university}} \\
\midrule
Takes A-level economics               &              & 0.152\sym{***}  & 0.158\sym{***}    \\
                                      &              & (0.010)         & (0.010)           \\
\addlinespace
A-level economics available            & 0.009\sym{***} &              &                \\
                                      & (0.0006)         &              &                \\
\addlinespace
Observations                          & 417,845           & 417,851           & 419,207             \\
\midrule
\multicolumn{4}{c}{\textbf{Panel B: Effect on log earnings at age 28, by degree pathway}} \\
\midrule
                                      & All          & Econ         & No econ        \\
                                      &              & degree       & degree         \\
                                      & (1)          & (2)          & (3)            \\
\midrule
Takes A-level economics               & 0.135\sym{*}  & -0.037  & 0.128             \\
                                      & (0.069)         & (0.056)         & (0.090)           \\
\addlinespace
Observations                          & 368,262           & 5,813           & 362,449             \\
Share of compliers                    &              & 15.2           & 84.8             \\
\bottomrule
\end{tabular}
\begin{minipage}{\textwidth}
\smallskip
\footnotesize
\textit{Notes:} Panel~A reports the effect of A-level economics take-up on the probability of studying economics at university. Column~1 is the reduced-form effect of availability; columns~2--3 instrument take-up with availability using the linear and saturated specifications. Panel~B reports the IV estimate of the effect of taking A-level economics on log real earnings at age~28, for the full sample (column~1, reproducing the linear IV estimate from Table~\ref{tab:main}) and separately for students who do and do not go on to study economics at university (columns~2--3). The split in Panel~B conditions on a post-treatment variable; the subgroup estimates are descriptive of where complier earnings gains are concentrated, not separately identified causal effects. Among compliers, 15.2\% shift into an economics degree (Panel~A) and 84.8\% do not. The subsample estimates need not aggregate exactly to the overall estimate, because each is estimated on a different sample; for the same reason the IV estimate in Panel~A differs slightly from the reduced form divided by the first stage. All specifications include school and cohort fixed effects and individual controls. Standard errors clustered at the school level in parentheses. Sample: students in switcher schools, 2002--2008 GCSE cohorts. \sym{*} $p<0.10$, \sym{**} $p<0.05$, \sym{***} $p<0.01$.
\end{minipage}
\end{table}
